\documentclass[xcolor=dvipsnames, compress, 12pt, aspectratio=169, presentation]{beamer}
\usepackage{../preamble_slides}
\newcommand{\sym}[1]{\ensuremath{^{\text{#1}}}}

\title{Staying Power: The Individual Effects of the H-1B Visa Lottery}
\author{Amy Kim \\ \small Princeton University, Department of Economics}
\date{Midwest Economics Association: SOLE \\ March 22, 2026}

\begin{document}
\begin{frame}[plain]
    \titlepage
    \addtocounter{framenumber}{-1}
\end{frame}

\section{Introduction}
\begin{frame}
    \frametitle{Motivation}
    \begin{itemize}
        \item Large literature: effects of immigration on U.S. labor market \textcolor{gray}{\footnotesize [Card 1990; Borjas 2003; Ottaviano and Peri 2012]}
        % \begin{itemize}
        %     \item Complicated by endogeneity of immigration
        % \end{itemize}
        % very extensive, very old literature in labor economics on the effects of immigration on the labor market
        \vspace{4mm}
        \item Common setting: \textbf{H-1B visa program} %common setting to study this question: h-1b visa lottery; three key features: 1) high skilled workers, 2) employer sponsorship, 3) lottery allocation; these features make it a compelling setting to study the effects of immigration
        \begin{itemize}
            \item Tied to employer-worker pair; allocated \textit{via lottery}
            \item Recent concerns about fraud + rising demand $\rightarrow$ new \$100k fee
        \end{itemize}
        \vspace{4mm}
        \item Existing literature: GE, local labor market, firm effects \textcolor{gray}{\footnotesize [Bound et al.\ 2017; Doran et al.\ 2022; Mahajan et al.\ 2024]} %conflicting results -- why? also, when we see firm-level responses to h1b lottery, mechanism hard to understand -- very important from policy perspective; at worker level, understanding what happens to workers who lose the lottery (CF) also important . fundamentally, for policymaking we want to know how these visas are actively being used by firms, and what the counterfactual is.
        \begin{itemize}
            \item How are firms using the lottery in practice? 
            \item What happens to workers who lose the lottery?
        \end{itemize}
        \vspace{4mm}
        \item \textcolor{maincolor}{\textbf{This Project:}} New application-level data + linkage strategy $\rightarrow$ first to estimate \textbf{worker-level effects} of winning H-1B lottery
    \end{itemize}
\end{frame}

\begin{frame}
    \frametitle{This Paper}
    \textbf{Research Question:} What are the effects of winning the H-1B lottery on the individual worker?
    \vspace{4mm}
    \begin{enumerate}
        \item Novel linked panel data: focus on \textbf{international students} applying for H-1B
        \only<1>{
        \begin{itemize}
            \item Large subset of modern H-1B applicants (40--50\%) are recent U.S. graduates
            \item Linked panel of 167k H-1B lottery applications to LinkedIn work + location histories
        \end{itemize}
        }
        \vspace{3mm}
        \item Estimate effects of lottery win on indiv. outcomes
        \only<2>{
        \begin{itemize}
            \item \textbf{Preview:} Winning $\uparrow$ U.S. retention by 1.52 pp (1.7\%) 2 yrs post-lottery
            \item \textcolor{red}{fill in with more results}
        \end{itemize}
        }
        \vspace{3mm}
        \item Investigating robustness and heterogeneity: why are effects small?
        \only<3>{
        \begin{enumerate}
            \item \textbf{Measurement Error:} Estimated effects 3-4x larger in higher-confidence linked samples, but still small
            \item \textbf{Student visa:} Effects largest for workers near student work authorization expiration, but still small
        \end{enumerate}
        }
    \end{enumerate} \vspace{4mm}
    \only<4>{\textbf{Empirical Puzzle:} How are lottery losers managing to stay?}
\end{frame}

\begin{frame}
    \frametitle{Literature Review}
    \begin{itemize}
        \item \textbf{Effects of H-1B Visa}
        \begin{itemize}
            \item Aggregate and firm-level impacts on native employment and wages \textcolor{gray}{\footnotesize [Kerr et al. 2015; Peri et al. 2015a; Doran et al. 2022; Mahajan et al. 2024]}, innovation and productivity \textcolor{gray}{\footnotesize [Kerr and Lincoln 2010; Peri et al. 2015b; Dimmock et al. 2022; Mayda et al. 2023]}, and employer monopsony power \textcolor{gray}{\footnotesize [Kim and Pei 2022]}
            \item Global effects on offshoring \textcolor{gray}{\footnotesize [Glennon 2024]}, origin countries \textcolor{gray}{\footnotesize [Khanna and Morales 2023]}, and third-party countries \textcolor{gray}{\footnotesize [Brinatti and Guo 2024]} 
            \item \textbf{This paper:} first individual-level evidence
        \end{itemize}
        \vspace{4mm}
        \item \textbf{International Students and U.S. Labor Markets}
        \begin{itemize}
            \item 
            \item Offshoring, Indian IT industry, immigration to third countries \textcolor{gray}{\footnotesize [Glennon 2024; Khanna and Morales 2023; Brinatti and Guo 2024]}
            \item Brain drain and return migration \textcolor{gray}{\footnotesize [Amanzadeh et al.\ 2024; Adda et al.\ 2022]}
            \item \textbf{This paper:} quantifies effect of H-1B cap on individual workers globally
        \end{itemize}
        \vspace{4mm}
        \item \textbf{Probabilistic Record Linkage} \textcolor{gray}{\footnotesize [Fellegi and Sunter 1969; Bailey et al.\ 2020; Enamorado et al.\ 2019]}
        \begin{itemize}
            \item \textbf{This paper:} new method for deidentified immigration records
        \end{itemize}
    \end{itemize}
\end{frame}


%%% TABLE OF CONTENTS
\begin{frame}[plain]{Outline}
    \tableofcontents
    \addtocounter{framenumber}{-1}
\end{frame}

%%% TABLE OF CONTENTS -- CURRENT SECTION
\AtBeginSection[ ]
{
\begin{frame}[plain]{Outline}
    \tableofcontents[currentsection, hideothersubsections]
    \addtocounter{framenumber}{-1}
\end{frame}
}

%%% TABLE OF CONTENTS -- CURRENT SECTION
\AtBeginSubsection[ ]
{
\begin{frame}[plain]{Outline}
    \tableofcontents[currentsubsection]
    \addtocounter{framenumber}{-1}
\end{frame}
}

\section{Background and Context}
\begin{frame}
    \frametitle{What is an H-1B Visa?}
    \begin{itemize}
        \item Temporary work visa (6 years) for workers in `specialty occupations'
        \begin{itemize}
            \item Requires highly specialized knowledge + at least a bachelor's degree
        \end{itemize}
        \vspace{3mm}
        \item Employer-sponsored: firm must apply on worker's behalf
        \begin{itemize}
            \item Portable between employers, but no self-sponsorship
            \item Employer can sponsor worker for green card (but application tied to employer)
        \end{itemize}
        \vspace{3mm}
        \item Capped at 85,000 visas per year since 2004
        \begin{itemize}
            \item 65,000 regular cap + 20,000 advanced degree exempt (ADE) for U.S.\ M.A.+
            \item Demand has exceeded cap every year since 2013
        \end{itemize}
        \vspace{3mm}
        \item Since 2013: visas allocated via \textbf{lottery} each spring
        \begin{itemize}
            \item Results in April, visa valid starting October
        \end{itemize}
    \end{itemize}
\end{frame}

\begin{frame}
    \frametitle{The H-1B Lottery: Key Facts}
    \begin{itemize}
        \item FY2021: USCIS moved to electronic registration (\$10 fee vs.\ \$500+ full petition)
        \begin{itemize}
            \item \textcolor{maincolor}{\textbf{Critical for this paper:}} USCIS began keeping records of lottery \textit{losers}
            \item Prior data: only winners observed
        \end{itemize}
        \vspace{3mm}
        \item Electronic system enabled fraud: duplicate applications via different employers
        \begin{itemize}
            \item Outsourcing/staffing firms `overapply' --- treat workers as interchangeable
            \item I exclude firm-years with any duplicate applicants
        \end{itemize}
        \vspace{3mm}
        \item Win rates: $\sim$35--40\% for regular lottery
        \begin{itemize}
            \item ADE applicants entered into separate lottery first (higher win rate)
        \end{itemize}
        \vspace{3mm}
        \item \textbf{Identification:} within firm-year, lottery selection is as good as random
    \end{itemize}
\end{frame}

\begin{frame}
    \frametitle{Optional Practical Training: The Key Bridge}
    \begin{itemize}
        \item F-1 (student) visa holders may work after graduation under OPT
        \begin{itemize}
            \item Standard OPT: 12 months of authorized employment
            \item STEM extension: up to 36 months total for STEM graduates
        \end{itemize}
        \vspace{3mm}
        \item OPT has become dominant pathway into H-1B
        \begin{itemize}
            \item 50--70\% of recent H-1B approvals go to workers already in U.S.
            \item 94\% of F-1$\rightarrow$H-1B transitions use OPT
        \end{itemize}
        \vspace{3mm}
        \item \textbf{Common trajectory:} graduate $\rightarrow$ OPT $\rightarrow$ apply for H-1B each spring (up to 3 tries)
        \begin{itemize}
            \item Worker already employed at sponsoring firm during OPT
            \item Firm retains worker if they win; worker faces uncertainty if they lose
        \end{itemize}
        \vspace{3mm}
        \item \textcolor{maincolor}{\textbf{Key insight:}} pre-lottery employment relationship enables linkage
    \end{itemize}
\end{frame}

\section{Data and Linkage}
\begin{frame}
    \frametitle{Data Sources}
    \begin{itemize}
        \item \textbf{H-1B Application Data} (FOIA via Bloomberg, FY2020--2023)
        \begin{itemize}
            \item All applicants: country of birth, birth year, sex, employer name/location
            \item Winners only: current visa type, education level/field, job title/salary
            \item 1.8 million applications; first dataset with both winners and losers
        \end{itemize}
        \vspace{4mm}
        \item \textbf{Revelio Labs LinkedIn Data}
        \begin{itemize}
            \item Profiles (894m individuals): full name, education history, work history
            \item Positions (1.7b records): employer, title, start/end dates, location, imputed salary
            \item Cross-sectional snapshot; not directly identified with H-1B status
        \end{itemize}
        \vspace{4mm}
        \item \textbf{Challenge:} H-1B data is deidentified --- no names, no unique identifiers
        \begin{itemize}
            \item Standard text matching infeasible
        \end{itemize}
    \end{itemize}
\end{frame}

\begin{frame}
    \frametitle{Linkage: Using Pre-Lottery Employment as an Anchor}
    \textbf{Key insight:} restrict candidates to workers already at sponsoring firm before lottery
    \vspace{3mm}
    \begin{enumerate}
        \item \textbf{Employer crosswalk:} fuzzy match H-1B firms to Revelio firms (93\% match rate)
        \begin{itemize}
            \item Exclude firm-years with 10+ applications or any duplicate applicants
        \end{itemize}
        \vspace{2mm}
        \item \textbf{Revelio imputation:} estimate country of birth, birth year, sex, graduation year
        \begin{itemize}
            \item Country: from education history location (high school) + name-based model
            \item Birth year: inferred from education timing
            \item Sex: name-based probabilistic imputation
        \end{itemize}
        \vspace{2mm}
        \item \textbf{Candidate scoring:} for each applicant, score LinkedIn workers at same firm
        \begin{itemize}
            \item Weights: country of birth (70\%), birth year (20\%), sex (10\%)
            \item Convert scores to within-applicant weights summing to 1
        \end{itemize}
    \end{enumerate}
    \vspace{2mm}
    \textbf{Result:} weighted panel of 167,598 linked applications
\end{frame}

\begin{frame}
    \frametitle{Sample Variants}
    \begin{itemize}
        \item \textbf{One-to-one} \textcolor{maincolor}{(preferred):} optimal bipartite matching (N = 64,781)
        \begin{itemize}
            \item Each applicant assigned to at most one Revelio candidate
        \end{itemize}
        \vspace{3mm}
        \item \textbf{Baseline:} 167,598 linked applications
        \begin{itemize}
            \item Restrict to U.S.-educated applicants (not U.S. high school); 70,580 with firm-year FE
        \end{itemize}
        \vspace{3mm}
        \item \textbf{Rare Country:} exclude high-multiplicity countries (N = 10,440)
        \begin{itemize}
            \item Drop applicants from Southern Asia, East Asia, UK, Canada, Australia
            \item Lower candidate pool multiplicity $\rightarrow$ sharper individual identification
        \end{itemize}
        \vspace{3mm}
        \item \textbf{Strict (Q25 / Q50 / Q75):} score threshold samples
        \begin{itemize}
            \item Restrict to applicants whose best match score exceeds the 25th, 50th, or 75th percentile
            \item Higher confidence matches; trade off sample size for precision
        \end{itemize}
    \end{itemize}
\end{frame}

\begin{frame}
    \frametitle{Sample Variants: Match Validation}
    \vfill
    \begin{center}
        Higher-confidence samples show better agreement between FOIA fields and matched Revelio profiles.\\[6mm]
        \hyperlink{match_quality}{\beamerbutton{Match Quality Table}}
    \end{center}
    \vfill
\end{frame}

\section{Results}
\begin{frame}
    \frametitle{Empirical Specification}
    For applicant $i$ at firm $c$ in lottery year $t$, estimate:
    $$y_{ict}^s = \beta_0 + \beta_1 Win_{ict} + \beta_2 ADE_{it} + \gamma_{ct} + \varepsilon_{ict}^s$$
    \vspace{2mm}
    \begin{itemize}
        \item $Win_{ict}$: lottery selection indicator
        \vspace{2mm}
        \item $ADE_{it}$: advanced degree exempt (separate, higher-odds lottery)
        \vspace{2mm}
        \item $\gamma_{ct}$: firm-by-lottery-year fixed effects
        \vspace{2mm}
        \item Outcomes measured $s$ years post-lottery; weights from probabilistic linkage
        \vspace{2mm}
        \item Standard errors clustered by firm
    \end{itemize}
    \vspace{3mm}
    \textbf{Identification:} within firm-year, lottery selection is quasi-random conditional on ADE status
\end{frame}

\begin{frame}
    \hypertarget{winning_lottery}{}
    \frametitle<1>{Winning the Lottery Stabilizes U.S. Presence}
    \frametitle<2>{Winning the Lottery Stabilizes U.S. Presence and Firm Retention}
    \frametitle<3>{Winning the Lottery Stabilizes U.S. Presence, Firm Retention, and Education}
    \vfill
    \begin{center}
    \resizebox{\textwidth}{!}{%
    \begin{tabular}{lcccccc}
    \toprule
    & In U.S. & In home ctry & Non-US non-home & Same firm & New pos.\ (cond.) & New educ. \\
    & (1) & (2) & (3) & (4) & (5) & (6) \\
    \midrule
    winner & 0.0116\sym{***} & $-$0.0060\sym{***} & $-$0.0033\sym{***} & \visible<2->{0.0264\sym{***}} & \visible<2->{$-$0.0146\sym{*}} & \visible<3->{$-$0.0042} \\
     & (0.0027) & (0.0015) & (0.0011) & \visible<2->{(0.0053)} & \visible<2->{(0.0075)} & \visible<3->{(0.0059)} \\
    \midrule
    \textbf{Control mean} & 0.950 & 0.019 & 0.010 & \visible<2->{0.743} & \visible<2->{0.448} & \visible<3->{0.340} \\
    \textbf{N} & 64{,}781 & 64{,}781 & 64{,}781 & \visible<2->{64{,}781} & \visible<2->{45{,}209} & \visible<3->{64{,}781} \\
    \bottomrule
    \end{tabular}
    }
    \end{center}
    \begin{center}
    \scriptsize\textit{OLS estimates; one-to-one linked sample (N = 64,781); firm-by-lottery-year FE; SEs clustered by firm. $t = 1$ year post-lottery. $^*p<0.10$, $^{**}p<0.05$, $^{***}p<0.01$.}\\[2mm]
    \hyperlink{appendix_additional_outcomes}{\beamerbutton{Additional outcomes}}
    \end{center}
    \vfill
\end{frame}

\begin{frame}
    \hypertarget{eventtime_in_us}{}
    \frametitle{Effect on U.S. Presence Grows as OPT Runway Shrinks}
    \vfill
    \begin{center}
        \includegraphics[height=0.73\textheight]{/home/yk0581/output/reg/graphs/event_time_us_educ_opt_in_us.png}\\[2mm]
        \hyperlink{rawmeans_in_us}{\beamerbutton{Raw means}}
    \end{center}
    \vfill
\end{frame}

\begin{frame}
    \hypertarget{eventtime_still_at_firm}{}
    \frametitle{Winners Remain Locked Into Sponsoring Firm Over Time}
    \vfill
    \begin{center}
        \includegraphics[height=0.73\textheight]{/home/yk0581/output/reg/graphs/event_time_us_educ_opt_still_at_firm.png}\\[2mm]
        \hyperlink{rawmeans_still_at_firm}{\beamerbutton{Raw means}}
    \end{center}
    \vfill
\end{frame}

%%% TAKING STOCK PART 1: REASON 1
\begin{frame}
    \frametitle{Taking Stock: Why Are Effects Small?}
    \begin{itemize}
        \item Effects are statistically significant but moderate in the one-to-one sample (1.52 pp U.S.\ presence)
        \vspace{4mm}
        \item \textbf{Reason 1:} Measurement error attenuates estimates
        \begin{itemize}
            \item Probabilistic linkage mixes correct and incorrect matches
            \item \textbf{How do effects vary by linkage quality?}
        \end{itemize}
        \vspace{4mm}
        \item \textbf{Reason 2: OPT workers have fallback options}
    \end{itemize}
\end{frame}

%%% ROBUSTNESS TABLE
\begin{frame}
    \hypertarget{robustness_table}{}
    \frametitle{Effects Robust Across Higher-Confidence Linked Samples}
    \vfill
    \begin{center}
    \resizebox{0.75\textwidth}{!}{%
    \begin{tabular}{lcccc}
    \toprule
    & In U.S. & Same firm & New pos.\ (cond.) & New educ. \\
    & (1) & (2) & (3) & (4) \\
    \midrule
    % One-to-one (overlay 1)
    \visible<1->{One-to-one ($N=64{,}781$)} & \visible<1->{0.0116\sym{***} {\scriptsize(0.0027)}} & \visible<1->{0.0264\sym{***} {\scriptsize(0.0053)}} & \visible<1->{$-$0.0146\sym{*} {\scriptsize(0.0075)}} & \visible<1->{$-$0.0042 {\scriptsize(0.0059)}} \\[4pt]
    % Baseline (overlay 2)
    \visible<2->{Baseline ($N=70{,}580$)} & \visible<2->{0.0023\sym{***} {\scriptsize(0.0008)}} & \visible<2->{0.0058\sym{***} {\scriptsize(0.0017)}} & \visible<2->{$-$0.0010 {\scriptsize(0.0020)}} & \visible<2->{0.0040\sym{**} {\scriptsize(0.0020)}} \\[4pt]
    % Strict Q25--Q75 (overlay 3)
    \visible<3->{Strict Q25 ($N=40{,}165$)} & \visible<3->{0.0103\sym{***} {\scriptsize(0.0026)}} & \visible<3->{0.0253\sym{***} {\scriptsize(0.0048)}} & \visible<3->{$-$0.0062 {\scriptsize(0.0063)}} & \visible<3->{$-$0.0073 {\scriptsize(0.0059)}} \\[4pt]
    \visible<3->{Strict Q50 ($N=27{,}765$)} & \visible<3->{0.0146\sym{***} {\scriptsize(0.0034)}} & \visible<3->{0.0328\sym{***} {\scriptsize(0.0060)}} & \visible<3->{$-$0.0169\sym{**} {\scriptsize(0.0081)}} & \visible<3->{$-$0.0085 {\scriptsize(0.0075)}} \\[4pt]
    \visible<3->{Strict Q75 ($N=13{,}894$)} & \visible<3->{0.0142\sym{***} {\scriptsize(0.0049)}} & \visible<3->{0.0408\sym{***} {\scriptsize(0.0086)}} & \visible<3->{$-$0.0336\sym{***} {\scriptsize(0.0119)}} & \visible<3->{0.0041 {\scriptsize(0.0113)}} \\[4pt]
    % Rare Country (overlay 4)
    \visible<4->{Rare Country ($N=10{,}440$)} & \visible<4->{0.0083\sym{**} {\scriptsize(0.0033)}} & \visible<4->{0.0221\sym{***} {\scriptsize(0.0061)}} & \visible<4->{$-$0.0097 {\scriptsize(0.0072)}} & \visible<4->{$-$0.0087 {\scriptsize(0.0069)}} \\
    \bottomrule
    \end{tabular}
    }
    \end{center}
    \begin{center}
    \scriptsize\textit{$\hat\beta_1$ (winner effect) with SEs in parentheses; firm-by-lottery-year FE; SEs clustered by firm. $t = 1$ year post-lottery. $^*p<0.10$, $^{**}p<0.05$, $^{***}p<0.01$.}\\[2mm]
    \hyperlink{robustness_ctrlmeans}{\beamerbutton{Control means}}
    \end{center}
    \vfill
\end{frame}

%%% TAKING STOCK PART 2: REASON 2
\begin{frame}
    \frametitle{Taking Stock: Why Are Effects Small?}
    \begin{itemize}
        \item Effects are statistically significant but moderate in the one-to-one sample (1.16 pp U.S.\ presence)
        \vspace{4mm}
        \item \textbf{Reason 1: Measurement error attenuates estimates}
        \vspace{4mm}
        \item \textbf{Reason 2: OPT workers have fallback options}
        \begin{itemize}
            \item Selecting on OPT $\implies$ applicants can continue to work through OPT!
            \item \textbf{Prediction:} effects largest when applicants have no OPT remaining
            \begin{itemize}
                \item Workers 1 yr post-graduation: nearing end of initial 12-month OPT
                \item Workers 3 yrs post-graduation: nearing end of STEM-OPT extension
            \end{itemize}
            \item Effects larger for firms less likely to re-sponsor workers (\textit{no-rep}) vs.\ \textit{high-rep} firms
            \item \textbf{How do effects vary by years since graduation and employer type?}
        \end{itemize}
    \end{itemize}
\end{frame}

%%% HETEROGENEITY TABLE
\begin{frame}
    \frametitle{Effects Strongest for Workers Near OPT Expiration}
    \vfill
    \begin{center}
    \resizebox{\textwidth}{!}{%
    \begin{tabular}{lcccccc}
    \toprule
    & 0 yrs post-grad & 1 yr post-grad & 2 yrs post-grad & 3 yrs post-grad & No rep. & High rep. \\
    & (1) & (2) & (3) & (4) & (5) & (6) \\
    \midrule
    winner & 0.0210 & 0.0324\sym{**} & 0.0395\sym{**} & 0.0446\sym{**} & 0.0221 & 0.0029 \\
     & (0.0195) & (0.0138) & (0.0164) & (0.0194) & (0.0211) & (0.0064) \\
    \midrule
    \textbf{Control mean} & 0.893 & 0.880 & 0.870 & 0.875 & 0.889 & 0.920 \\
    \textbf{N} & 1{,}980 & 6{,}029 & 4{,}728 & 3{,}043 & 2{,}605 & 19{,}322 \\
    \bottomrule
    \end{tabular}
    }
    \end{center}
    \begin{center}
    \scriptsize\textit{Outcome: \textit{In U.S.}\ (2 years post-lottery). Cols (1)--(4): subgroups by years since graduation at time of lottery application. Cols (5)--(6): by employer repetition. $^*p<0.10$, $^{**}p<0.05$, $^{***}p<0.01$.}
    \end{center}
    \vfill
\end{frame}

\begin{frame}
    \frametitle{Summary of Results}
    \begin{itemize}
        \item \textbf{Main effects:} winning H-1B raises U.S.\ presence (+1.52 pp) and firm retention (+1.23 pp)
        \begin{itemize}
            \item Winners less likely to leave for a different employer ($-$1.23 pp)
            \item Winners less likely to appear in home country ($-$0.74 pp)
        \end{itemize}
        \vspace{4mm}
        \item \textbf{Robustness:} effects consistent across higher-confidence samples
        \begin{itemize}
            \item U.S.\ presence reaches 2.0 pp (Q75); firm retention 1.45 pp
            \item Baseline effects smaller --- consistent with attenuation bias from probabilistic linkage
        \end{itemize}
        \vspace{4mm}
        \item \textbf{Heterogeneity:} effects driven by workers near OPT expiration
        \begin{itemize}
            \item Patterns at 2 and 4 yrs post-graduation consistent with OPT expiration dates
        \end{itemize}
        \vspace{4mm}
        \item \textbf{Interpretation:} H-1B as \textit{stability lottery} --- reduces disruption, locks workers into firm
    \end{itemize}
\end{frame}

\section{Conclusion}
\begin{frame}
    \frametitle{Conclusion}
    \begin{itemize}
        \item \textbf{Contribution 1:} first individual-level causal estimates of H-1B lottery effects
        \begin{itemize}
            \item Winning stabilizes U.S.\ presence and firm retention; reduces employer mobility
            \item For OPT workers, H-1B is a stability mechanism, not an entry gate
        \end{itemize}
        \vspace{3mm}
        \item \textbf{Contribution 2:} novel probabilistic linkage for deidentified immigration records
        \begin{itemize}
            \item Pre-lottery employment as anchor reduces candidate space from millions to hundreds
        \end{itemize}
        \vspace{3mm}
        \item \textbf{Policy implication:} modest average effects $\neq$ unimportant policy
        \begin{itemize}
            \item Repeated reapplication costs, employer dependence, and uncertainty have welfare costs
            \item Reforms targeting portability and OPT-to-H1B transitions could reduce disruption
        \end{itemize}
        \vspace{3mm}
        \item \textbf{Next steps:}
        \begin{itemize}
            \item Wage and compensation effects
            \item Productivity outcomes via GitHub data
            \item Return migration and third-country reallocation
        \end{itemize}
    \end{itemize}
\end{frame}

%%% BACKUP SLIDES

\begin{frame}[noframenumbering]
    \hypertarget{appendix_additional_outcomes}{}
    \frametitle{Additional Outcomes: One-to-One Sample}
    \vfill
    \begin{center}
    \resizebox{\textwidth}{!}{%
    \begin{tabular}{lcccc}
    \toprule
    & Profile in U.S. & Cont.\ education & Activity update diff & Frac.\ null end date \\
    & (1) & (2) & (3) & (4) \\
    \midrule
    winner & \textcolor{gray}{---} & 0.0040 & 0.7620\sym{***} & $-$0.0129\sym{**} \\
     & \textcolor{gray}{---} & (0.0059) & (0.2922) & (0.0056) \\
    \midrule
    \textbf{Control mean} & \textcolor{gray}{---} & 0.663 & 17.011 & 0.408 \\
    \textbf{N} & \textcolor{gray}{---} & 64{,}781 & 64{,}781 & 64{,}781 \\
    \bottomrule
    \end{tabular}
    }
    \end{center}
    \begin{center}
    \scriptsize\textit{OLS estimates; one-to-one linked sample (N = 64,781); firm-by-lottery-year FE; SEs clustered by firm. $t = 1$ year post-lottery. $^*p<0.10$, $^{**}p<0.05$, $^{***}p<0.01$.}\\[2mm]
    \hyperlink{winning_lottery}{\beamerbutton{Return}}
    \end{center}
    \vfill
\end{frame}

\begin{frame}[noframenumbering]
    \hypertarget{robustness_ctrlmeans}{}
    \frametitle{Robustness: Control Means by Sample}
    \vfill
    \begin{center}
    \resizebox{0.75\textwidth}{!}{%
    \begin{tabular}{lcccc}
    \toprule
    & In U.S. & Same firm & New pos.\ (cond.) & New educ. \\
    \midrule
    One-to-one   & 0.950 & 0.743 & 0.448 & 0.340 \\
    Baseline     & 0.956 & 0.743 & 0.388 & 0.301 \\
    Strict Q25   & 0.938 & 0.741 & 0.442 & 0.452 \\
    Strict Q50   & 0.934 & 0.740 & 0.449 & 0.466 \\
    Strict Q75   & 0.937 & 0.746 & 0.455 & 0.473 \\
    Rare Country & 0.948 & 0.749 & 0.413 & 0.304 \\
    \bottomrule
    \end{tabular}
    }
    \end{center}
    \begin{center}
    \scriptsize\textit{Control group means; $t = 1$ year post-lottery.}\\[2mm]
    \hyperlink{robustness_table}{\beamerbutton{Return}}
    \end{center}
    \vfill
\end{frame}

\begin{frame}[noframenumbering]
    \hypertarget{match_quality}{}
    \frametitle{Higher-Quality Samples Have Better Match Validation}
    \vfill
    \begin{center}
    \resizebox{0.85\textwidth}{!}{%
    \input{/home/yk0581/output/reg/tables/paper_match_quality.tex}
    }
    \end{center}
    \begin{center}
    \scriptsize\textit{Agreement rates between FOIA petition fields and matched Revelio profiles, by sample.}
    \end{center}
    \vfill
\end{frame}

\begin{frame}[noframenumbering]
    \hypertarget{rawmeans_in_us}{}
    \frametitle{Raw Means: U.S. Presence by Lottery Outcome}
    \vfill
    \begin{center}
        \includegraphics[height=0.73\textheight]{/home/yk0581/output/reg/graphs/event_time_raw_us_educ_opt_in_us.png}\\[2mm]
        \hyperlink{eventtime_in_us}{\beamerbutton{Return}}
    \end{center}
    \vfill
\end{frame}

\begin{frame}[noframenumbering]
    \hypertarget{rawmeans_still_at_firm}{}
    \frametitle{Raw Means: Firm Retention by Lottery Outcome}
    \vfill
    \begin{center}
        \includegraphics[height=0.73\textheight]{/home/yk0581/output/reg/graphs/event_time_raw_us_educ_opt_still_at_firm.png}\\[2mm]
        \hyperlink{eventtime_still_at_firm}{\beamerbutton{Return}}
    \end{center}
    \vfill
\end{frame}

\end{document}
